% Partie : méthodes
% Chapitre : démonstrations
% Section : employer les définitions

\documentclass[../../../include/open-logic-section]{subfiles}

\begin{document}

\olfileid[fr]{mth}{prf}{def}

\olsection{Employer les définitions}

Nous avons dit qu'il faut connaître
toutes les définitions susceptibles
d'intervenir dans une démonstration
et savoir les appliquer correctement.
Ce point est essentiel ; examinons-le
de plus près. Les définitions servent
à abréger des propriétés et des
relations pour que nous puissions
en parler plus succinctement.
L'expression introduite est le
\emph{definiendum} ; ce qu'elle
abrège est le \emph{definiens}.
Dans une démonstration, nous devons
souvent revenir à la manière dont
le definiendum a été introduit :
pour avancer, il faut exploiter
la structure logique du definiens,
c'est-à-dire de la version longue
abrégée par le terme défini.
Déplier les définitions permet
ainsi d'atteindre le cœur du
raisonnement logique.

Commençons par un exemple.
Supposons que vous vouliez
démontrer l'énoncé suivant :

\begin{prop}
Pour tous ensembles $A$ et $B$,
$A \cup B = B \cup A$.
\end{prop}

Pour commencer, examinons le signe «~$=$~» dans cette égalité.
L'extensionnalité donne un critère d'égalité des ensembles.
Deux ensembles sont identiques si et seulement s'ils ont les
mêmes !!{element}s. Nous utiliserons ce critère sous la forme
suivante :

\begin{defn}
Les ensembles $A$ et $B$ sont
\emph{identiques}, $A = B$,
si et seulement si tout !!{element}
de~$A$ est !!a{element} de~$B$,
et réciproquement.
\end{defn}

Dans ce critère, $A$ et~$B$ désignent des ensembles quelconques.
L'expression «~$A = B$~» joue le rôle du \emph{definiendum}.
Nous la remplaçons par la condition suivante :
«~tout !!{element} de~$A$
est !!a{element} de~$B$,
et réciproquement~».
Cette condition est le
\emph{definiens}.\footnote{
Dans ce cas particulier — et
c'est assez déroutant —, si
$A = B$, les ensembles $A$
et $B$ ne sont qu'un seul
et même ensemble, bien que
nous lui donnions des lettres
différentes à gauche et à
droite. Mais les manières de
désigner cet ensemble peuvent
différer, ce qui rend la
définition non triviale.}
La définition dit que
$A = B$ est vraie si et
seulement si la condition
est remplie. Nous abrégeons
«~si et seulement si~»
par «~ssi~».

Pour appliquer la définition,
il faut faire correspondre
ses variables $A$ et $B$
aux expressions du cas traité.
Ici, pour que
$A \cup B = B \cup A$
soit vraie, chaque
$z \in A \cup B$ doit
aussi appartenir à
$B \cup A$, et
réciproquement. Dans
la proposition,
$A \cup B$ joue le rôle
du~$A$ de la définition,
et $B \cup A$ celui
du~$B$. Puisque les mêmes
lettres $A$ et $B$
figurent dans la définition
et dans la proposition,
mais avec des fonctions
différentes, il faut veiller
à ne pas les confondre.
Ce serait, par exemple,
une erreur de penser qu'il
suffit de montrer que
tout !!{element} de~$A$
est !!a{element} de~$B$,
et réciproquement :
cela démontrerait $A = B$,
et non
$A \cup B = B \cup A$.
Comme $A$ et $B$ peuvent
être deux ensembles quelconques,
on ne pourrait d'ailleurs
pas aller loin sans hypothèse
supplémentaire : ils peuvent
très bien être différents.

Notre démonstration utilise
également des notions
ensemblistes comme l'union :
il nous faut donc connaître
le sens du symbole $\cup$
pour savoir comment procéder.
Il arrive que le dépliage
d'une définition en appelle
d'autres. Ainsi,
$A \cup B$ est défini
comme
$\Setabs{z}{z \in A
\text{ ou } z \in B}$.
Pour démontrer que
$x \in A \cup B$,
le dépliage de la définition
de $\cup$ montre qu'il
faut établir
$x \in \Setabs{z}{z \in A
\text{ ou } z \in B}$.
Il faut alors se rappeler
que
$x \in \Setabs{z}{\dots z\dots}$
si et seulement si
$\dots x\dots$.
En dépliant aussi la
notation
$\Setabs{z}{\dots z \dots}$,
nous voyons que ce qu'il
faut montrer est
$x \in A$ ou $x \in B$.
Ainsi, «~tout !!{element}
de $A \cup B$ est aussi
!!a{element} de $B \cup A$~»
signifie en réalité :
«~pour tout $x$, si
$x \in A$ ou $x \in B$,
alors $x \in B$ ou
$x \in A$~».
Si nous déplions toutes
les définitions de la
proposition, voici ce
qu'il faut démontrer :

\begin{prop}
Pour tous ensembles $A$
et $B$ : (a) pour tout
$x$, si $x \in A$ ou
$x \in B$, alors
$x \in B$ ou
$x \in A$ ; et
(b) pour tout $x$,
si $x \in B$ ou
$x \in A$, alors
$x \in A$ ou
$x \in B$.
\end{prop}

L'essentiel est que déplier
les définitions constitue
une étape nécessaire dans
la construction d'une
démonstration. Ce n'est
pas toujours facile :
il faut distinguer avec
soin les variables de
la définition et les
termes de l'énoncé à
démontrer, puis les faire
correspondre. Pour réussir,
il faut comprendre la
question posée et le
sens de tous les termes
qu'elle emploie. Souvent,
il faut déplier plus
d'une définition.
Dans les démonstrations
simples qui suivent,
la solution découle
presque immédiatement
des définitions.
Ce ne sera évidemment
pas toujours aussi simple.

\begin{prob}
Supposons que l'on vous
demande de démontrer que
$A \cap B \neq \emptyset$.
Dépliez toutes les
définitions qui interviennent,
c'est-à-dire reformulez
cet énoncé sans employer
«~$\cap$~», «~$=$~»
ni «~$\emptyset$~».
\end{prob}

\end{document}
